% 图契约：RLC滑杆的统一方向与方程；图中的端口、几何、箭头和方程对应本题。
% 版面：沿用前三模块的几何图+推导分区；目标印宽17cm，简图至中等密度。
% 布局：左侧真实结构，右侧关键关系；真实电路边不作为架构流程箭头处理。
\documentclass[tikz,border=7pt]{standalone}
\usepackage{fontspec,xeCJK,amsmath,amssymb}
\IfFontExistsTF{Times New Roman}{\setmainfont{Times New Roman}}{\setmainfont{TeX Gyre Termes}}
\IfFontExistsTF{Heiti SC}{\setCJKmainfont[AutoFakeBold=1.5]{Heiti SC}}{\setCJKmainfont[AutoFakeBold=1.5]{FandolHei-Regular}}
\usetikzlibrary{arrows.meta,bending,calc,patterns,positioning,decorations.pathmorphing}
\definecolor{ink}{HTML}{243745}
\definecolor{blue}{HTML}{275DAB}
\definecolor{green}{HTML}{177C63}
\definecolor{red}{HTML}{BD4935}
\definecolor{orange}{HTML}{AA6C16}
\definecolor{purple}{HTML}{715393}
\tikzset{every node/.style={font=\fontsize{9.3}{12.5}\selectfont,text=ink},
arrow/.style={-{Stealth[length=5pt 1.5,width'=0pt 3]},shorten >=2pt,shorten <=1pt,line width=1pt},
arrow short/.style={-{Stealth[length=3pt 1.0,width'=0pt 2.5]},shorten >=1pt,shorten <=0pt,line width=.8pt},
axis/.style={arrow,draw=ink!60,line width=.6pt},
vec/.style={arrow,draw=blue},
wire/.style={draw=ink!75,line width=.8pt},
guide/.style={dashed,draw=ink!45,line width=.5pt},
card/.style={rounded corners=2pt,draw=blue!65,fill=blue!4,align=center,text width=6.5cm,inner xsep=9pt,inner ysep=7pt},
title/.style={font=\fontsize{10.5}{14}\selectfont\bfseries,text width=15cm,align=center},
label/.style={font=\fontsize{8.4}{11}\selectfont,align=center},
note/.style={label,text width=6.1cm}}
\begin{document}
\begin{tikzpicture}[x=1cm,y=1cm]
\path[use as bounding box] (-8,-3.5) rectangle(8,3.65);
\node[title] at (0,3.25) {理论图\quad RLC滑杆的统一方向与方程};
\begin{scope}[shift={(-4.05,0)}]\draw[wire](-2,1.5)--(1.3,1.5)(-2,-1.5)--(1.3,-1.5);\draw[ink,line width=2pt](1.3,-1.5)--(1.3,1.5);\draw[wire,draw=ink](-2,1.5)--(-2.0,1.21) (-2.0,0.79)--(-2,0.5);\begin{scope}[shift={(-2.0,1.0)},rotate=-90.0]\draw[draw=ink,fill=white,line width=.8pt](-0.21,-.085)rectangle(0.21,.085);\end{scope}
\node[label] at(-1.72,1.0){$R$};
\begin{scope}[shift={(-2,0.5)},rotate=-90]\draw[wire](0,0)--(0.3,0);\draw[wire](0.3,0)..controls(0.3,0.5)and(0.4,0.5)..(0.4,0);\draw[wire](0.4,0)..controls(0.4,0.5)and(0.5,0.5)..(0.5,0);\draw[wire](0.5,0)..controls(0.5,0.5)and(0.6000000000000001,0.5)..(0.6000000000000001,0);\draw[wire](0.6000000000000001,0)..controls(0.6000000000000001,0.5)and(0.7,0.5)..(0.7,0);\draw[wire](0.7,0)--(1.0,0);\node[below]at(0.5,-0.2){$$};\end{scope}\node[left]at(-2.4,0.0){$L$};\draw[wire](-2,-0.5)--(-2,-0.9)(-2,-1.1)--(-2,-1.5);\draw[ink,thick](-2.45,-0.9)--(-1.55,-0.9)(-2.45,-1.1)--(-1.55,-1.1);\node[left]at(-2.5,-1.0){$C$};\draw[vec,green](1.3,2)--(2.8,2)node[midway,above]{$v$};\draw[vec,red](1.3,-.6)--(1.3,.65)node[right]{$i$};\draw[vec,orange](1,.2)--(-.4,.2)node[midway,above]{$F_B$};\node[blue]at(0,-.8){$\otimes B$};\draw[<->,orange](2.5,-1.5)--(2.5,1.5)node[midway,right]{$\ell$};\end{scope}\node[card,draw=blue!65,fill=blue!4]at(3.9,1.8){$L\dot i+Ri+q/C=Gv$};\node[card,draw=green!65,fill=green!4]at(3.9,0){$m\dot v=F_{\rm ext}-Gi$};\node[card,draw=orange!65,fill=orange!4]at(3.9,-1.8){$G=B\ell,\quad \dot q=i$};
\end{tikzpicture}\end{document}